% ===================================================================
%  TABELL Nr 3 — Barndom och ungdom.
%  Traduction 2026 d'apres texto/io, controlee sur texto/fr et
%  texto/es. Consigne : voir texto/sv/10-tabell-01.tex.
% ===================================================================

%%K t03-apar-1 apar
\VUsaut{3.0mm}
\VUtitre{15.0pt}{60}{TABELL Nr 3}
\VUsaut{1.6mm}
\VUtitre{11.4pt}{0}{Barndom och ungdom.}
\VUsaut{3.4mm}

%%K t03-apar-2 apar
(Den 3\textsuperscript{e} och 4\textsuperscript{e} tabellen är så sammansatta att de i fyra
\VUgras{familjescener} framställer de väsentliga begreppen om
\VUgras{vädret}, \VUgras{åldern}, \VUgras{familjen}, \VUgras{klädseln} och
\VUgras{ståndet}.)
\VUsaut{4.0mm}

%%K t03-tit-1 sub
\VUcentre{12.0pt}{0}{\textit{Första scenen.}}
\VUsaut{3.0mm}

%%K t03-tit-2 sub
\VUpk{10.2pt}{40}{Dopceremoni i en by.}
\VUsaut{2.4mm}

%%K t03-tit-3 sub
\VUpk{10.2pt}{40}{Våren.}
\VUsaut{3.0mm}

%%K t03-c1-01-1 p
1. --- Vi är i \VUgras{våren}, i månaderna \VUgras{mars}, \VUgras{april} eller
\VUgras{maj}, på \VUgras{veckans} dagar, \VUgras{måndag}, \VUgras{tisdag},
eller \VUgras{onsdag}. Klockan är \VUgras{tio} på morgonen.

%%K t03-c1-02-1 p
2. --- \VUgras{Vädret} är vackert, \VUgras{luften} ren. Solen skiner. Några
\VUgras{små moln} \textsuperscript{(2)} stiger upp på den blå \VUgras{himlen}
\textsuperscript{(1)}.

%%K t03-c1-03-1 p
3. --- \VUgras{Träden} har knappt några \VUgras{löv}. De smäckra
\VUgras{popplarna} \textsuperscript{(3)} och den höga \VUgras{platanen}
\textsuperscript{(17)} har hittills bara knoppar.

%%K t03-c1-04-1 p
4. --- \VUgras{Svalorna} \textsuperscript{(4)} kommer tillbaka och flyger med glatt kvitter
kring \VUgras{spiran} \textsuperscript{(7)} på \VUgras{klocktornet} \textsuperscript{(5)} som
kröner byns \VUgras{kyrka} \textsuperscript{(6)}.

%%K t03-tit-4 sub
\VUsaut{3.4mm}
\VUpk{10.2pt}{40}{Kyrkan.}
\VUsaut{3.0mm}

%%K t03-c1-05-1 p
5. --- Mycket enkel är den kyrkan med sin stora \VUgras{portal} \textsuperscript{(13)} och
tjocka \VUgras{ur} \textsuperscript{(8)}. Den lilla
\VUgras{visaren} \textsuperscript{(9)} står på siffran X, den visar
\VUgras{timmarna}. Den \VUgras{stora} \textsuperscript{(10)} står på XII
(middag); den visar \VUgras{minuterna}.

%%K t03-c1-06-1 p
6. --- I klocktornet ringer \VUgras{klockorna} \textsuperscript{(11)} glatt. På takets topp
står ett litet \VUgras{kors} \textsuperscript{(12)} av sten.

%%K t03-c1-07-1 p
7. --- Kyrkan har inte bara en stor \VUgras{portal} \textsuperscript{(13)}, den har också en
liten sidodörr. Genom den lilla dörren kommer herr
\VUgras{kyrkoherden} \textsuperscript{(15)} klädd i sin \VUgras{kaftan} ut ur
\VUgras{sakristian} \textsuperscript{(14)} och går till \VUgras{prästgården}
\textsuperscript{(19)}.

%%K t03-c1-08-1 p
8. --- Bredvid honom, här är \VUgras{mjölkerskan} \textsuperscript{(24)} med sin
\VUgras{åsna} \textsuperscript{(23)}. Hon kommer tillbaka från staden, dit hon
förde god mjölk åt barnen.

%%K t03-tit-5 sub
\VUsaut{4.0mm}
\VUpk{10.2pt}{40}{Fruktträdgården.}
\VUsaut{3.4mm}

%%K t03-c1-09-1 p
9. --- Herr Aliboron (åsnan) är helt nöjd med den morgonturen. Han andas med förtjusning in
luften som är balsamerad av doften från \VUgras{blommorna} \textsuperscript{(20)} som täcker
\VUgras{fruktträden} \textsuperscript{(18)} i den närliggande \VUgras{fruktträdgården}. Helt
rosa och vita är de
\VUgras{persikoträden} \textsuperscript{(18)} och \VUgras{aprikosträden}
\textsuperscript{(19)}, och de låter hoppas på en god skörd.

%%K t03-c1-10-1 p
10. --- Under tiden går de små \VUgras{bina} \textsuperscript{(22)} från
\VUgras{bikupan} \textsuperscript{(21)} dit för att surrande röva sin söta
\VUgras{honung}.

%%K t03-c1-11-1 p
11. --- Men vad händer? Här kommer byns barn springande och skrämmer bort de förskräckta
\VUgras{gässen} \textsuperscript{(25)}. Det är följets utgång ur kyrkan efter \VUgras{dopet}
av notariens son.

%%K t03-c1-12-1 p
12. --- Lille Jakob vaknar i sin \VUgras{vagga} \textsuperscript{(26)} av klockornas glada
ringning, och hans syster Magdalena som också vill se dopföljet, ber sin mor komma för att
kamma hennes hår och klä henne på husets tröskel.

%%K t03-c1-13-1 p
13. --- Modern samtycker. Hon sätter på sig sin \VUgras{huvudduk} \textsuperscript{(28)}, sitt
\VUgras{liv} \textsuperscript{(29)} och sitt
\VUgras{förkläde} \textsuperscript{(30)}, och hon kommer och sätter sig på en
liten stol framför dörren. Magdalena har bara sin \VUgras{underkjol} \textsuperscript{(31)}
och sin \VUgras{korsett} \textsuperscript{(32)}.

%%K t03-c1-14-1 p
14. --- Så lycklig hon är! Hon glömmer sina älsklingsblommor, sina
\VUgras{nejlikor} \textsuperscript{(34)}, på vilka \VUgras{fjärilarna}
\textsuperscript{(35)} slår sig ner, sina \VUgras{tulpaner} \textsuperscript{(36)}, sina
\VUgras{liljekonvaljer} \textsuperscript{(37)}, i
\VUgras{krukor} \textsuperscript{(38)}, på \VUgras{muren}
\textsuperscript{(33)}, och sina \VUgras{rosor} \textsuperscript{(39)} som bildar en så vacker
häck. Hon betraktar de rika kläderna hos damerna i följet.

%%K t03-c1-15-1 p
15. --- Byns barn ger inte akt så mycket på kläderna. De rusar fram och knuffar varandra för
att få \VUgras{karameller} \textsuperscript{(55)} eller
\VUgras{mynt} \textsuperscript{(52)}. En av dem gråter \textsuperscript{(40)}.

%%K t03-c1-16-1 p
16. --- Den andre trampade på tassen på \VUgras{hunden} \textsuperscript{(42)} som skrikande
flyr och skrämmer bort \VUgras{hönan} \textsuperscript{(43)} och hennes \VUgras{kycklingar}
\textsuperscript{(44)}.

%%K t03-tit-6 sub
\VUsaut{5.0mm}
\VUpk{10.2pt}{40}{Dopföljet.}
\VUsaut{4.0mm}

%%K t03-c1-17-1 p
17. --- Framför följet går \VUgras{amman} \textsuperscript{(45)}. Hon har en vacker
\VUgras{hätta} \textsuperscript{(46)} med långa band. Hon bär
\VUgras{den nyfödde} \textsuperscript{(47)} helt klädd i \VUgras{spetsar}
\textsuperscript{(48)}.

%%K t03-c1-18-1 p
18. --- Bakom amman kommer \VUgras{gudfadern} \textsuperscript{(49)} och
\VUgras{gudmodern} \textsuperscript{(53)}. De delar ut karamellerna och mynten.
Gudfadern har \VUgras{handskar} \textsuperscript{(51)} och en \VUgras{cylinderhatt}
\textsuperscript{(50)}. Gudmodern håller en vacker \VUgras{bukett} \textsuperscript{(54)} i
handen.

%%K t03-c1-19-1 p
19. --- Bakom dem ser vi barnets \VUgras{farbror} \textsuperscript{(56)} och
\VUgras{faster} \textsuperscript{(58)}. Fastern har en elegant \VUgras{hatt}
\textsuperscript{(59)}, farbrodern en svart \VUgras{bonjour} \textsuperscript{(57)} och en vit
väst. Här är sedan \VUgras{kusinen} \textsuperscript{(60)} och \VUgras{kusinen}
\textsuperscript{(61)}. Denna sistnämnda håller vid handen \VUgras{systern}
\textsuperscript{(62)} till den nyfödde, lilla Berta.

%%K t03-c1-20-1 p
20. --- Bertas andre \VUgras{bror} \textsuperscript{(63)} går bakom. Han ger handen åt en
\VUgras{släkting} \textsuperscript{(64)}. Familjens
\VUgras{vänner} \textsuperscript{(65)} kommer sedan.

%%K t03-tit-7 sub
\VUsaut{4.6mm}
\VUpk{10.2pt}{40}{Åskådarna.}
\VUsaut{3.4mm}

%%K t03-c1-21-1 p
21. --- Bredvid följet, här är en \VUgras{tiggare} \textsuperscript{(66)} som stöder sig på
sin \VUgras{krycka} \textsuperscript{(67)}. Han ber om en allmosa.

%%K t03-c1-22-1 p
22. --- På andra sidan står en liten mycket \VUgras{artig} \textsuperscript{(68)} pojke. Han
hälsar och önskar « \VUgras{god dag} » åt de personer han känner.

%%K t03-c1-23-1 p
23. --- Byborna har gått ut ur sina hem för att se dopföljet. Vi ser en
\VUgras{bonde} \textsuperscript{(69)} med sin \VUgras{bomullsmössa} och bredvid
en \VUgras{bondkvinna} \textsuperscript{(70)}. Sedan en \VUgras{gammal kvinna}
\textsuperscript{(71)} som stöder sig på sin \VUgras{käpp} \textsuperscript{(72)}. Slutligen
\VUgras{mjölnaren} \textsuperscript{(73)} med sin stora \VUgras{filthatt}
\textsuperscript{(74)} och en ung bondkvinna med
\VUgras{träskor} \textsuperscript{(75)} på fötterna, som lär sitt lilla barn att
gå.

%%K t03-c1-24-1 p
24. --- Den scenen är mycket rolig.
\VUsaut{4.0mm}

%%K t03-tit-8 sub
\VUcentre{12.0pt}{0}{\textit{Andra scenen.}}
\VUsaut{3.0mm}

%%K t03-tit-9 sub
\VUpk{10.2pt}{40}{Offentlig fest.}
\VUsaut{2.4mm}

%%K t03-tit-10 sub
\VUpk{10.2pt}{40}{Vattenlekar och kapprodd.}
\VUsaut{3.0mm}

%%K t03-c2-01-1 p
1. --- \VUgras{Sommaren} har kommit. Den heta solen i månaden \VUgras{juni} (juli eller
\VUgras{augusti}) eggar de unga att bada och ro.

%%K t03-c2-02-1 p
2. --- På flodens högra strand klär roddarna av sig för att delta i vattenlekarna och
kapprodden.

%%K t03-c2-03-1 p
3. --- En av dem tar av sig sina \VUgras{skor} \textsuperscript{(1)}; han knyter upp (knäpper
upp) dem. Hans två kamrater är redan färdiga. Nu har de bara sin \VUgras{tröja}
\textsuperscript{(2)} och \VUgras{badbyxa} \textsuperscript{(3)}. De samtalar medan de väntar
på sina vänner. De talar om marknaden och festen som äger rum på andra stranden.

%%K t03-c2-04-1 p
4. --- På marken, bredvid dem, ligger kamraternas \VUgras{kläder}: ett
\VUgras{par} \VUgras{kängor} \textsuperscript{(4)}, \VUgras{skor}
\textsuperscript{(5)}, \VUgras{strumpor} \textsuperscript{(6)},
\VUgras{filthattar} \textsuperscript{(7)} och \VUgras{halmhattar}
\textsuperscript{(8)} och en \VUgras{halsduk} \textsuperscript{(9)}.

%%K t03-c2-05-1 p
5. --- En av de unga har omsorgsfullt vikt ihop sin \VUgras{rock} \textsuperscript{(10)}. Han
lägger den på en handduk, i vilken hans granne snart ska svepa in de två \VUgras{plaggen}.

%%K t03-c2-06-1 p
6. --- En sjätte pojke är sen. Han har fortfarande sin \VUgras{mössa} \textsuperscript{(11)}
på huvudet. Han har varken tagit av sig sin
\VUgras{skjorta} \textsuperscript{(12)}, eller sin \VUgras{krage}
\textsuperscript{(13)}, eller sina \VUgras{manschetter} \textsuperscript{(14)}. Han håller sin
\VUgras{väst} \textsuperscript{(17)} i handen och har fortfarande sina \VUgras{byxor}
\textsuperscript{(16)} och \VUgras{hängslen} \textsuperscript{(15)}. Han kommer säkert inte
att vara färdig till nästa tävling.

%%K t03-c2-07-1 p
7. --- Bredvid de unga leder en stig med många \VUgras{tistlar} \textsuperscript{(18)} på
sidorna ner till flodens strand, där fader Jakob mellan
\VUgras{vassen} \textsuperscript{(19)} gör i ordning \VUgras{fyrverkeriet}
\textsuperscript{(20)} som man ska skjuta av på kvällen.

%%K t03-tit-11 sub
\VUsaut{4.4mm}
\VUpk{10.2pt}{40}{Tävlingen.}
\VUsaut{3.4mm}

%%K t03-c2-08-1 p
8. --- På floden gör några damer, som skyddar sig med sina parasoller, en tur i en
\VUgras{båt} \textsuperscript{(21)}. De följer tävlingen som är mycket spännande. I varje
\VUgras{jolle} \textsuperscript{(22)} ror fyra
\VUgras{roddare} \textsuperscript{(24)} av alla krafter, medan \VUgras{styrmannen}
\textsuperscript{(26)} håller \VUgras{rodret} \textsuperscript{(25)} och styr den bräckliga
roddbåten. \VUgras{Årorna} \textsuperscript{(23)} slår i takt mot vattnet och de lätta båtarna
far fram med svindlande fart.

%%K t03-c2-09-1 p
9. --- Några unga kämpar med varandra, bredvid brons pelare, om priset på
\VUgras{den insmorda masten} \textsuperscript{(28)}. En av dem går försiktigt på
den böjliga och hala masten för att nå den lilla \VUgras{flaggan} \textsuperscript{(29)} som
är fäst vid änden. Hans olycklige \VUgras{medtävlare} \textsuperscript{(30)} har fallit i
vattnet. Han simmar för att komma i land.

%%K t03-c2-10-1 p
10. --- Äldre unga håller \VUgras{en (båt)tornering} \textsuperscript{(32)} vid
\VUgras{kajen} \textsuperscript{(33)}. Alla dessa lekar bevistas av en talrik
\VUgras{publik} \textsuperscript{(31)} som lutar sig mot \VUgras{räcket} på
\VUgras{bron} \textsuperscript{(27)} och trängs på \VUgras{stranden}. Några
åskådare vänder sig dock om för att följa leken med \VUgras{prismasten} \textsuperscript{(45)}
eller för att beundra \VUgras{stadsstyrelsens följe} som kommer för \VUgras{prisutdelningen}.

%%K t03-c2-11-1 p
11. --- Först, här är några \VUgras{gatpojkar} \textsuperscript{(35)} som går före
\VUgras{musikanterna} \textsuperscript{(36)}; sedan kommer
\VUgras{borgmästaren} \textsuperscript{(37)}, hans ställföreträdare och den
lilla stadens \VUgras{stadsfullmäktige}. Slutligen går \VUgras{brandmännen}
\textsuperscript{(38)}.

%%K t03-tit-12 sub
\VUsaut{5.0mm}
\VUpk{10.2pt}{40}{Marknaden.}
\VUsaut{3.6mm}

%%K t03-c2-12-1 p
12. --- Talrika människor är på \VUgras{marknadsplatsen} \textsuperscript{(39)}. Man kommer
för att se de olika \VUgras{stånden}: \VUgras{trähästarna} \textsuperscript{(40)},
\VUgras{cirkusen} \textsuperscript{(41)} där föreställningen redan har börjat;
\VUgras{den offentliga balen} \textsuperscript{(42)}, där stadens unga män och unga kvinnor
njuter av
\VUgras{dansens} nöje.

%%K t03-c2-13-1 p
13. --- Längst in ser vi \VUgras{arenan} \textsuperscript{(44)} där den tappre spanske
\VUgras{matadoren} kämpar mot den rasande \VUgras{tjuren} \textsuperscript{(45)}, sedan
\VUgras{rullskridskobanan} \textsuperscript{(49)},
\VUgras{skjutbanan} \textsuperscript{(46)} där man övar sig i att bruka
\VUgras{gevären} och att skjuta mot \VUgras{måltavlan}.

%%K t03-c2-14-1 p
14. --- Nära intill, här är \VUgras{marknadsteatern} \textsuperscript{(47)} i vilken man
spelar \VUgras{komedi} och \VUgras{drama}. Alldeles intill står
\VUgras{jättens} \textsuperscript{(48)} och \VUgras{dvärgens} stånd.
\VUgras{Längden} på den förste är två meter och femton centimeter; den andre är
knappt så stor som ett barn.

%%K t03-c2-15-1 p
15. --- Slutligen, här är det stora \VUgras{menageriet} \textsuperscript{(50)} med sina
\VUgras{vilda djur}, sin \VUgras{tiger} \textsuperscript{(51)} med kraftiga \VUgras{klor},
sitt \VUgras{lejon} \textsuperscript{(52)} med tät
\VUgras{man}, sina två \VUgras{giraffer} \textsuperscript{(53)} och sin
\VUgras{elefant} \textsuperscript{(54)} som så skickligt använder sin
\VUgras{snabel}.

%%K t03-c2-16-1 p
16. --- \VUgras{Djurtämjaren} \textsuperscript{(55)} som dresserar de djuren står vid ståndets
dörr.
\VUsaut{6.0mm}
\VUfilet{22mm}
